\subsection{Stokes Theory}

% **Role:** calculus as local linearization plus global geometric integration.
%
% Generalized Stokes theorem, Green's theorem, classical Stokes theorem, divergence theorem.

Stokes theory relates exterior differentiation to oriented boundary, with Green's theorem, classical Stokes, and divergence as special cases.

\begin{definition} \label{def:affine-chains-and-boundaries} \hlt{Chains and Boundaries}\\
For $\mathbf{p}_0,\ldots,\mathbf{p}_k\in\R^n$, oriented affine $k$-simplex $\sigma=[\mathbf{p}_0,\ldots,\mathbf{p}_k]$ is $k$-surface on $Q^k$: $\sigma(\mathbf{u})=\mathbf{p}_0+\sum\limits_{i=1}^{k}u_i(\mathbf{p}_i-\mathbf{p}_0)$.\\
Permuting vertices multiplies $\sigma$ by the permutation sign. An \hlt{affine $k$-chain} is a finite formal sum of oriented affine $k$-simplices; integration is extended linearly.\\
For $k\geq1$, define $\partial_{\mathrm{ch}}[\mathbf{p}_0,\ldots,\mathbf{p}_k] \defeq \sum\limits_{j=0}^{k}(-1)^j[\mathbf{p}_0,\ldots,\mathbf{p}_{j-1},\mathbf{p}_{j+1},\ldots,\mathbf{p}_k]$, and extend $\partial_{\mathrm{ch}}$ linearly. If $\Phi=T\circ\sigma$ is a $C^2$ oriented simplex, set $\partial_{\mathrm{ch}}\Phi\defeq T(\partial_{\mathrm{ch}}\sigma)$, again extended linearly to chains.
\end{definition}

\begin{theorem} \label{thm:stokes-theorem-differential-forms}\label{thm:forms-manifolds-stokes-main-source-todo} \hlt{Stokes' Theorem}\\
If $\Psi$ is a $C^2$ $k$-chain in open $V\subseteq\R^m$ and $\omega$ is a $C^1$ $(k-1)$-form in $V$, then $\int_{\Psi}d\omega=\int_{\partial_{\mathrm{ch}}\Psi}\omega$.
\begin{proof}
By linearity, take $\Psi=\Phi=T\circ\sigma$, where $\sigma$ is the standard oriented $k$-simplex.\\
By Theorem \ref{thm:pullback-basic-laws}, $\int_{\Phi}d\omega =\int_{\sigma}(d\omega)_T =\int_{\sigma}d(\omega_T)$; and $\int_{\partial_{\mathrm{ch}}\Phi}\omega =\int_{\partial_{\mathrm{ch}}\sigma}\omega_T$, remains to prove formula on $\sigma$.\\
By linearity, reduce to $\lambda=f\,dx_1\wedge\cdots\wedge\widehat{dx_r}\wedge\cdots\wedge dx_k$.\\
The boundary formula in Definition \ref{def:affine-chains-and-boundaries} leaves only the two faces in the $x_r$ direction, and the resulting identity is the one-variable fundamental theorem of calculus inside the iterated integral.\\
\hlt{Structural Commentary}: Pullback reduces Stokes' theorem to a computation on the standard simplex.\\
\hlt{Key Insight}: The boundary signs are exactly the signs produced by the missing wedge factor.
\end{proof}
\end{theorem}

\begin{definition} \label{def:closed-and-exact-forms} \hlt{Closed and Exact Forms}\\
A $k$-form $\omega$ in $E$ is \hlt{exact} if $\omega=d\lambda$ for some $(k-1)$-form $\lambda$ in $E$, and \hlt{closed} if $\omega$ is $C^1$ with $d\omega=0$.\\
By Theorem \ref{thm:exterior-derivative-basic-laws}, exact $C^1$ forms are closed.
\end{definition}

\begin{perspective} \label{per:closed-exact-local-global} \hlt{Closed Is Local, Exact Is Global}\\
Exact $\Rightarrow$ closed is just $d^2=0$; the converse depends on the domain.\\
On $\R^2\setminus\{\ze\}$ the standard warning example is $\omega=\frac{-y}{x^2+y^2}\,dx+\frac{x}{x^2+y^2}\,dy$, which is closed but not exact.\\
Holes can obstruct global antiderivatives even when local derivative tests vanish.
\end{perspective}

\begin{theorem} \label{thm:poincare-lemma-convex-rudin} \hlt{Poincare Lemma on Convex Sets}\\
If $E\subseteq\R^n$ is convex and open, then every closed $C^1$ $k$-form in $E$, $k\geq1$, is exact.\\
The same holds on every open set $C^2$-equivalent to a convex open set.
\begin{proof}
Choose $\mathbf{a}\in E$ and let $H_t(\x)=\mathbf{a}+t(\x-\mathbf{a})$, $0\leq t\leq1$. Convexity gives $H_t(E)\subseteq E$. For a $k$-form $\omega$, define the $(k-1)$-form  as $(K\omega)_{\x}(\mathbf{v}_1,\ldots,\mathbf{v}_{k-1}) \defeq \int_0^1 t^{k-1}\omega_{H_t(\x)}(\x-\mathbf{a},\mathbf{v}_1,\ldots,\mathbf{v}_{k-1})\,dt$.\\
A coordinate differentiation under the integral gives the homotopy identity $\omega-H_0^\ast\omega=dK\omega+Kd\omega$. Since $k\geq1$, the constant map $H_0$ pulls every $k$-form back to $0$. If $d\omega=0$, then $\omega=d(K\omega)$, so $\omega$ is exact.\\
If $E$ is $C^2$-equivalent to a convex open set $C$ by $\Phi:C\to E$, then $\Phi^\ast\omega$ is closed by Theorem \ref{thm:pullback-basic-laws}. The convex case gives $\Phi^\ast\omega=d\eta$ on $C$. Applying $(\Phi^{-1})^\ast$ and using the same pullback law gives $\omega=d((\Phi^{-1})^\ast\eta)$ on $E$.\\
\hlt{Structural Commentary}: Convexity lets coefficient-level antiderivatives be chosen consistently.\\
\hlt{Key Insight}: Closed is local symmetry; convexity supplies the global primitive.
\end{proof}
\end{theorem}

\begin{theorem} \label{thm:vector-analysis-forms-dictionary} \hlt{Vector Analysis Dictionary in $\R^3$}\\
For $\mathbf{F}=F_1\mathbf{e}_1+F_2\mathbf{e}_2+F_3\mathbf{e}_3$ on open $E\subseteq\R^3$,\\
set $\lambda_{\mathbf{F}}\defeq F_1\,dx+F_2\,dy+F_3\,dz$ and $\omega_{\mathbf{F}}\defeq F_1\,dy\wedge dz+F_2\,dz\wedge dx+F_3\,dx\wedge dy$. Then\\
$\mathbf{F}=\nabla u \Longleftrightarrow \lambda_{\mathbf{F}}=du$; and $\nabla\times\mathbf{F}=0 \Longleftrightarrow d\lambda_{\mathbf{F}}=0$; and\\
$\mathbf{F}=\nabla\times\mathbf{G} \Longleftrightarrow \omega_{\mathbf{F}}=d\lambda_{\mathbf{G}}$; and $\nabla\cdot\mathbf{F}=0 \Longleftrightarrow d\omega_{\mathbf{F}}=0$.\\
Hence $\nabla\times(\nabla u)=0$ and $\nabla\cdot(\nabla\times\mathbf{G})=0$; on sets $C^2$-equivalent to convex open sets, the converses hold.
\begin{proof}
Expand $d\lambda_{\mathbf{F}}$ and $d\omega_{\mathbf{F}}$ by Definition \ref{def:wedge-product-and-exterior-derivative}; their coefficients are exactly $\nabla\times\mathbf{F}$ and $\nabla\cdot\mathbf{F}$. Potentials and vector potentials are exactness of $\lambda_{\mathbf{F}}$ and $\omega_{\mathbf{F}}$; the converses follow from Theorem \ref{thm:poincare-lemma-convex-rudin}.\\
\hlt{Structural Commentary}: Gradient, curl, and divergence are exterior derivative in degrees $0,1,2$.\\
\hlt{Key Insight}: The familiar vector identities are $d^2=0$ written in $\R^3$ notation.
\end{proof}
\end{theorem}

\begin{theorem} \label{thm:green-stokes-divergence-special-cases} \hlt{Green, Classical Stokes, and Divergence}\\
The following are special cases of Theorem \ref{thm:stokes-theorem-differential-forms}.
\begin{enumerate}[label=(\roman*)]
\setlength{\itemsep}{0pt}
\item If $\Omega\subseteq\R^2$ has positively oriented boundary, $\alpha,\beta\in C^1$ near $\Omega$, then $\int_{\operatorname{bd}\Omega}(\alpha\,dx+\beta\,dy) = \int_{\Omega}\left(\frac{\partial\beta}{\partial x}-\frac{\partial\alpha}{\partial y}\right)\,dA$
\item If $\Phi$ is a $C^2$ $2$-surface in $\R^3$ and $\mathbf{F}\in C^1$ near $\Phi$, then $\int_{\Phi}(\nabla\times\mathbf{F})\cdot\mathbf{n}\,dA = \int_{\partial_{\mathrm{ch}}\Phi}\mathbf{F}\cdot\mathbf{t}\,ds$.
\item If $\Omega\subseteq\R^3$ has positively oriented boundary and $\mathbf{F}\in C^1$ near $\Omega$, then $\int_{\Omega}(\nabla\cdot\mathbf{F})\,dV =\int_{\operatorname{bd}\Omega}\mathbf{F}\cdot\mathbf{n}\,dA$
\end{enumerate}
\begin{proof}
\begin{enumerate}[label=(\roman*)]
\setlength{\itemsep}{0pt}
\item Put $\lambda=\alpha\,dx+\beta\,dy$. Then $d\lambda=(\partial\beta/\partial x-\partial\alpha/\partial y)\,dA$.
\item Put $\lambda=\lambda_{\mathbf{F}}$. By Theorem \ref{thm:vector-analysis-forms-dictionary}, $d\lambda=\omega_{\nabla\times\mathbf{F}}$.
\item Put $\omega=\omega_{\mathbf{F}}$. Then $d\omega=(\nabla\cdot\mathbf{F})\,dx\wedge dy\wedge dz$.
\end{enumerate}
Theorem \ref{thm:stokes-theorem-differential-forms}, together with the line, area, and volume element translations, gives the three displayed formulas.\\
\hlt{Structural Commentary}: Forms translate Stokes into the classical vector formulas.\\
\hlt{Key Insight}: Boundary plus $d$ is the content; vector notation is translation.
\end{proof}
\end{theorem}

\begin{method} \label{meth:forms-manifolds-proof-patterns-source-todo}\label{meth:differential-forms-proof-patterns} \hlt{Differential-Forms Proof Patterns}
\begin{enumerate}[label=(\roman*)]
\setlength{\itemsep}{0pt}
\item \hlt{Signs}: compute wedge signs by determinant swaps and permutation parity.
\item \hlt{Standard presentation}: remove repeated indices and sort the rest.
\item \hlt{Pullback}: substitute $dt_i$ for $dy_i$ and apply the chain rule.
\item \hlt{Stokes}: pull back to the standard simplex, prove the coordinate case, extend linearly.
\item \hlt{Closed versus exact}: use $d^2=0$ for necessity; use convexity or $C^2$-equivalence for primitives.
\item \hlt{Vector analysis}: translate fields to forms, apply Stokes, translate back.
\end{enumerate}
\hlt{Key Insight}: Forms package signs, orientation, and parameter changes.
\end{method}
